% ==============================================================================
% M.Tech Project (MTP) Comprehensive Viva & Defense Q&A Compendium
% Project: Sedimentation Dynamics of Rigid Multiblobs in Stokes Flow
% Document: Master's Defense Potential Questions, Mathematical Derivations & Model Answers
% ==============================================================================

\documentclass[11pt,a4paper]{article}
\usepackage[utf8]{inputenc}
\usepackage{amsmath,amssymb,amsfonts}
\usepackage{booktabs}
\usepackage{siunitx}
\usepackage{hyperref}
\usepackage{geometry}
\usepackage{xcolor}
\geometry{margin=1in}

\hypersetup{
    colorlinks=true,
    linkcolor=blue!80!black,
    citecolor=blue!80!black,
    urlcolor=blue!80!black
}

\newcommand{\question}[2]{
    \vspace{0.4cm}
    \noindent\colorbox{blue!10}{\parbox{\dimexpr\textwidth-2\fboxsep\relax}{
        \textbf{#1: #2}
    }}
    \vspace{0.2cm}
}

\newcommand{\shortanswer}[1]{
    \noindent\textbf{\textcolor{green!50!black}{[Key Defense Response]:}} #1
    \vspace{0.2cm}
}

\newcommand{\deepdetails}[1]{
    \noindent\textbf{\textcolor{blue!60!black}{[Mathematical \& Physical Derivation]:}} #1
    \vspace{0.2cm}
}

\newcommand{\pitfall}[1]{
    \noindent\textbf{\textcolor{red!70!black}{[Trap / What NOT to Say]:}} #1
    \vspace{0.3cm}
}

\title{\textbf{M.Tech Thesis Defense Compendium: Potential Questions, In-Depth Mathematical Answers, and Viva Strategies}}
\author{M.Tech Project Research Group \\ Department of Mechanical Engineering / Fluid Mechanics}
\date{November 2026}

\begin{document}
\maketitle

\begin{abstract}
This document serves as an exhaustive, advanced preparation guide for the Master's Thesis defense and viva-voce examination on \emph{Multi-Body Sedimentation Dynamics in Low Reynolds Number Regimes}. It compiles 41 high-probability technical, theoretical, and numerical questions frequently posed by examination committees and faculty panels. Questions span the continuum Navier-Stokes scaling analysis, Green's function formalism, regularized Rotne-Prager-Yamakawa (RPY) tensors, saddle-point GMRES formulations, the classical Crowley instability (1971), the celestial Keplerian orbit analogies discovered by Rahul Chajwa and co-workers (PRL 2019, PRX 2020), industrial and environmental use-cases (inkjet printing, oceanic carbon sequestration), and the extrapolation to machine learning surrogate modeling (Random Forest and XGBoost). For each question, a concise verbal response, rigorous mathematical justification, workspace simulation citations, and common examiner traps are provided.
\end{abstract}

\tableofcontents
\newpage

% ==============================================================================
\section{Category 1: Governing Equations, Low Reynolds Number \& Scaling Analysis}
% ==============================================================================

\question{Q1}{Why did you neglect the non-linear term $(\mathbf{u} \cdot \nabla)\mathbf{u}$ and unsteady acceleration $\partial \mathbf{u}/\partial t$ in the Navier-Stokes equations? Under what exact physical conditions is this justified?}

\shortanswer{We operate in the creeping flow regime where the Reynolds number is vanishingly small ($Re = \rho U L / \eta \ll 1$). Non-dimensionalizing the Navier-Stokes equations by characteristic length $L$, terminal velocity $U$, and viscous pressure scale $p_c = \eta U / L$ reveals that the inertial terms are multiplied by $Re$. In the asymptotic limit $Re \to 0$, inertial forces become completely negligible compared to viscous dissipation.}

\deepdetails{
Starting from the incompressible momentum equation:
\begin{equation}
    \rho \left( \frac{\partial \mathbf{u}}{\partial t} + (\mathbf{u} \cdot \nabla) \mathbf{u} \right) = -\nabla p + \eta \nabla^2 \mathbf{u} + \mathbf{f}_{\text{ext}}
\end{equation}
Introducing dimensionless variables $\mathbf{x}^* = \mathbf{x}/L$, $\mathbf{u}^* = \mathbf{u}/U$, $t^* = t U / L$, and $p^* = (p - p_\infty)/(\eta U / L)$, the equation transforms to:
\begin{equation}
    Re \left( \frac{\partial \mathbf{u}^*}{\partial t^*} + (\mathbf{u}^* \cdot \nabla^*) \mathbf{u}^* \right) = -\nabla^* p^* + \nabla^{*2} \mathbf{u}^* + \mathbf{f}^*_{\text{ext}}
\end{equation}
For a colloidal particle of radius $R \approx 10\,\mu\text{m}$ settling in water ($\eta = 10^{-3}\,\text{Pa}\cdot\text{s}$, $\rho = 10^3\,\text{kg/m}^3$) at $U \approx 10\,\mu\text{m/s}$:
\begin{equation}
    Re = \frac{10^3 \times 10^{-5} \times 10^{-5}}{10^{-3}} = 10^{-4} \ll 1
\end{equation}
Thus, inertial terms represent less than $0.01\%$ of viscous stresses, mathematically justifying their complete elimination.
}

\pitfall{Do not say ``fluid velocity is zero''. The fluid velocity is non-zero; it is the \emph{inertial acceleration} of fluid parcels that is negligible compared to the viscous shearing resistance.}

\vspace{0.4cm}
\question{Q2}{Why did you scale pressure by $p_c = \eta U / L$ rather than the traditional Bernoulli dynamic pressure $p_c = \frac{1}{2}\rho U^2$?}

\shortanswer{In high Reynolds number flows, pressure gradients balance inertial convection, making $\rho U^2$ the natural pressure scale. In low Reynolds number flows ($Re \ll 1$), pressure gradients directly balance viscous stresses ($\nabla p \sim \eta \nabla^2 \mathbf{u}$). Scaling pressure by $\rho U^2$ would lead to a singular perturbation where viscous terms blow up as $1/Re$.}

\deepdetails{
If one mistakenly scales pressure as $p^* = p / (\rho U^2)$, the momentum equation becomes:
\begin{equation}
    \frac{\partial \mathbf{u}^*}{\partial t^*} + (\mathbf{u}^* \cdot \nabla^*) \mathbf{u}^* = -\nabla^* p^* + \frac{1}{Re} \nabla^{*2} \mathbf{u}^*
\end{equation}
As $Re \to 0$, the viscous term $\frac{1}{Re} \nabla^{*2} \mathbf{u}^* \to \infty$, obscuring the dominant balance. The Stokes scaling $p_c = \eta U / L$ reflects the true physical balance between pressure and viscous shear stress ($\nabla p \approx \eta \nabla^2 \mathbf{u}$), keeping all non-inertial terms of order $\mathcal{O}(1)$.
}

\pitfall{Never say ``pressure doesn't matter in Stokes flow''. Pressure is a crucial Lagrange multiplier that enforces the incompressibility constraint $\nabla \cdot \mathbf{u} = 0$.}

\vspace{0.4cm}
\question{Q3}{You stated that Stokes flow is ``instantaneous'' and ``time-reversible''. What does that mean physically for a sedimenting particle? Does a particle dropped from rest accelerate?}

\shortanswer{Because the unsteady fluid term $\rho \partial \mathbf{u}/\partial t$ vanishes, the fluid has zero memory; flow fields adjust instantaneously to boundary motion. In a continuum fluid with $Re = 0$, a sedimenting particle reaches its steady terminal velocity instantaneously ($t = 0$) upon application of gravity. Kinematic time-reversibility means that if external forces are reversed, fluid streamlines and particle trajectories retrace themselves identically in reverse.}

\deepdetails{
In classical Newtonian mechanics, $m d\mathbf{v}/dt = \mathbf{F}_{\text{ext}} - \mathbf{F}_{\text{drag}}$. The particle momentum relaxation timescale is:
\begin{equation}
    \tau_p = \frac{m}{6\pi\eta R} = \frac{\frac{4}{3}\pi R^3 \rho_p}{6\pi\eta R} = \frac{2}{9} \frac{\rho_p R^2}{\eta}
\end{equation}
For a $1\,\mu\text{m}$ particle in water, $\tau_p \sim 10^{-7}\,\text{seconds}$. For all observable macroscopic integration timesteps ($\Delta t \ge 10^{-4}\,\text{s}$), the particle is strictly at terminal velocity ($d\mathbf{V}/dt \approx 0$). This explains why our quasi-steady mobility solve matches time-integration schemes to $10^{-14}$ precision.
}

\pitfall{Do not confuse fluid reversibility with thermodynamic reversibility. Viscous dissipation ($\Phi = 2\eta \mathbf{E}:\mathbf{E} > 0$) continually converts mechanical work into heat; energy is irreversibly dissipated, but kinematic streamlines are geometrically reversible under velocity inversion.}

% ==============================================================================
\section{Category 2: Green's Functions, Singularities \& Multiblob Discretization}
% ==============================================================================

\question{Q4}{What physically is a Green's function for Stokes flow, and why does it enable replacing a 3D volumetric mesh with boundary points?}

\shortanswer{A Green's function (the Stokeslet or Oseen tensor) is the exact analytical velocity field produced at any point in space by an infinitesimal point-force applied to the fluid. Because the Stokes equations are strictly linear, the Lorentz reciprocal theorem allows the flow field around any complex 3D body to be expressed as a boundary surface convolution of Green's functions with boundary surface tractions, reducing a 3D PDE into a 2D surface boundary integral.}

\deepdetails{
The governing singular problem is:
\begin{equation}
    \eta \nabla^2 \mathbf{u} - \nabla p = -\mathbf{F} \delta(\mathbf{x} - \mathbf{x}_0), \quad \nabla \cdot \mathbf{u} = 0
\end{equation}
The free-space solution is the Oseen tensor:
\begin{equation}
    \mathbf{u}(\mathbf{x}) = \mathbf{G}(\mathbf{x} - \mathbf{x}_0) \mathbf{F}, \quad \mathbf{G}(\mathbf{r}) = \frac{1}{8\pi\eta r} \left( \mathbf{I} + \frac{\mathbf{r}\mathbf{r}^T}{r^2} \right)
\end{equation}
By Green's representation formula:
\begin{equation}
    u_i(\mathbf{x}) = \int_{\partial \Omega} G_{ij}(\mathbf{x} - \mathbf{y}) \sigma_{jk}(\mathbf{y}) n_k \, dS_y - \int_{\partial \Omega} u_j(\mathbf{y}) T_{ijk}(\mathbf{x} - \mathbf{y}) n_k \, dS_y
\end{equation}
where $T_{ijk}$ is the stresslet tensor. For rigid bodies moving with velocity $\mathbf{U}$, this converts volumetric fluid solving into determining boundary tractions $\boldsymbol{\lambda}$.
}

\pitfall{Do not say Green's functions only work for spheres. Green's functions apply to \emph{arbitrary} geometries; boundary conditions dictate the surface traction distribution.}

\vspace{0.4cm}
\question{Q5}{The point Stokeslet diverges as $1/r$ when $r \to 0$. How does the Rotne-Prager-Yamakawa (RPY) tensor resolve this singularity, and why is it superior to naive point-force methods?}

\shortanswer{The RPY tensor regularizes the point singularity by treating each blob not as an infinitesimal mathematical point, but as a rigid spherical shell or Gaussian force distribution of finite hydrodynamic radius $a$. For overlapping blobs ($r < 2a$), the RPY formulation replaces the $1/r$ divergence with a bounded polynomial that preserves strict Symmetric Positive Definiteness (SPD), guaranteeing numerical stability.}

\deepdetails{
The RPY tensor $\mathbf{M}(\mathbf{r})$ is obtained by averaging the Stokeslet over the surfaces of two spheres of radius $a$ centered at separation $\mathbf{r}$:
\begin{equation}
    \mathbf{M}_{ij} = \frac{1}{(4\pi a^2)^2} \iint \mathbf{G}(\mathbf{r} + \mathbf{y}_i - \mathbf{y}_j) \, dS_i \, dS_j
\end{equation}
Evaluating this integral yields:
\begin{align}
    \mathbf{M}_{ij} &= \frac{1}{8\pi\eta r} \left[ \left(\mathbf{I} + \frac{\mathbf{r}\mathbf{r}^T}{r^2}\right) + \frac{2a^2}{r^2}\left(\frac{1}{3}\mathbf{I} - \frac{\mathbf{r}\mathbf{r}^T}{r^2}\right) \right], \quad (r \ge 2a) \\
    \mathbf{M}_{ij} &= \frac{1}{6\pi\eta a} \left[ \left(1 - \frac{9r}{32a}\right)\mathbf{I} + \frac{3}{32a}\frac{\mathbf{r}\mathbf{r}^T}{r} \right], \quad (r < 2a) \\
    \mathbf{M}_{ii} &= \frac{1}{6\pi\eta a} \mathbf{I}, \quad (r = 0)
\end{align}
Because $\mathbf{M}$ is SPD for all spatial configurations, the total fluid energy dissipation rate $\mathcal{P} = \boldsymbol{\lambda}^T \mathbf{M} \boldsymbol{\lambda} > 0$ for any non-zero force vector $\boldsymbol{\lambda}$, preventing unphysical negative drag or instability.
}

\pitfall{Never say RPY is just an empirical curve-fit. It is an exact analytical spatial convolution of the Stokeslet over spherical surfaces.}

\vspace{0.4cm}
\question{Q6}{How do you choose the blob radius $a$ for a given shell of geometric radius $R$ and $N$ blobs? What happens if $a$ is too small or too large?}

\shortanswer{We use the geometric area coverage formula: $a \approx \frac{1}{2}\sqrt{\frac{4\pi R_g^2}{N}}$. If $a$ is too small, the shell becomes hydrodynamically porous, allowing fluid to leak through the boundary and underpredicting drag. If $a$ is too large, excessive blob overlap artificially enlarges the effective hydrodynamic radius, overpredicting drag.}

\deepdetails{
To ensure that $N$ blobs of projected area $\pi a^2$ cover the spherical surface $4\pi R_g^2$ without substantial gaps or excessive artificial thickness, equating total blob cross-sectional area to surface area yields:
\begin{equation}
    N \cdot \pi a^2 \approx \pi R_g^2 \implies a \approx \frac{R_g}{\sqrt{N}} \implies a \approx \frac{1}{2} \sqrt{\frac{4\pi R_g^2}{N}}
\end{equation}
In our multi-resolution validation:
\begin{itemize}
    \item $N = 12 \implies a = 0.5117$, $R_h = 1.2625$ (porous shell limit, $V_0 = 42.36$)
    \item $N = 42 \implies a = 0.2735$, $R_h = 1.1220$ ($V_0 = 47.28$)
    \item $N = 162 \implies a = 0.1392$, $R_h = 1.0530$ ($V_0 = 50.35$)
    \item $N = 2562 \implies a = 0.0350$, $R_h = 1.0113$ (continuum limit, $V_0 = 52.45$, within $1.1\%$ of Stokes Law $53.05$).
\end{itemize}
}

% ==============================================================================
\section{Category 3: Numerical Integration Schemes \& GMRES Linear Solvers}
% ==============================================================================

\question{Q7}{In your presentation, you showed that Adams-Bashforth, Forward Euler, and the Direct Mobility scheme yielded identical terminal velocities to $10^{-14}$ precision. How is that possible when Adams-Bashforth is 2nd order and Euler is 1st order? Is that a bug?}

\shortanswer{No, it is a fundamental physical verification of steady Stokes flow. Because there is no fluid or particle inertia, a body under constant gravitational force settles at a constant velocity in unbounded fluid ($d\mathbf{V}/dt = 0$). For any ODE $\dot{y} = \text{const}$, 1st-order Euler, 2nd-order Adams-Bashforth, and direct algebraic inversion are mathematically identical.}

\deepdetails{
In unbounded Stokes flow with constant $\mathbf{F}_{\text{ext}}$, the rigid body equations are:
\begin{equation}
    \frac{d\mathbf{q}}{dt} = \mathbf{V}(\mathbf{q}) = \mathbf{N}(\mathbf{q}) \mathbf{F}_{\text{ext}}
\end{equation}
Because unbounded space is translationally invariant, the mobility tensor $\mathbf{N}$ does not depend on position $\mathbf{q}$. Therefore, $\mathbf{V} \equiv \mathbf{V}_0 = \text{constant}$.
\begin{itemize}
    \item Forward Euler: $\mathbf{V}^{n+1} = \mathbf{N} \mathbf{F} = \mathbf{V}_0$
    \item Adams-Bashforth: $\mathbf{V}^{n+1} = \frac{3}{2}\mathbf{V}^n - \frac{1}{2}\mathbf{V}^{n-1} = \frac{3}{2}\mathbf{V}_0 - \frac{1}{2}\mathbf{V}_0 = \mathbf{V}_0$
    \item Mobility Solve: $\mathbf{V} = (\mathbf{K}^T \mathbf{M}^{-1} \mathbf{K})^{-1} \mathbf{F}_{\text{ext}} = \mathbf{V}_0$
\end{itemize}
All schemes evaluate the exact same steady GMRES velocity. This proves our implementation introduces zero artificial temporal drift.
}

\pitfall{Do not say ``Adams-Bashforth is unnecessary''. It is necessary when particles sediment near a wall (where $\mathbf{N}(\mathbf{q})$ changes with height $z$) or in multi-body configurations where relative distances vary dynamically in time.}

\vspace{0.4cm}
\question{Q8}{What is the saddle-point system $\begin{pmatrix} \mathbf{M} & -\mathbf{K} \\ -\mathbf{K}^T & \mathbf{0} \end{pmatrix} \begin{pmatrix} \boldsymbol{\lambda} \\ \mathbf{U} \end{pmatrix} = \begin{pmatrix} \mathbf{0} \\ -\mathbf{F}_{\text{ext}} \end{pmatrix}$, and why do we use Preconditioned GMRES rather than direct matrix inversion?}

\shortanswer{The saddle-point system couples fluid mobility with rigid-body kinematic constraints. We use Preconditioned GMRES because direct matrix inversion scales as $\mathcal{O}((3N+6)^3)$, which becomes computationally prohibitive for $N > 100$. GMRES iteratively solves the system via matrix-vector products scaling as $\mathcal{O}(N^2)$ (or $\mathcal{O}(N \log N)$ with Fast Multipole Methods).}

\deepdetails{
The matrix size is $(3N + 6) \times (3N + 6)$:
\begin{itemize}
    \item $\mathbf{M} \in \mathbb{R}^{3N \times 3N}$ is the symmetric positive definite RPY mobility matrix;
    \item $\mathbf{K} \in \mathbb{R}^{3N \times 6}$ is the kinematic matrix where row $i$ is $\mathbf{K}_i = [\mathbf{I}, -(\mathbf{x}_i - \mathbf{q})_\times]$;
    \item $\boldsymbol{\lambda} \in \mathbb{R}^{3N}$ are the hydrodynamic constraint forces;
    \item $\mathbf{U} \in \mathbb{R}^6$ is the rigid body velocity vector $(\mathbf{V}^T, \mathbf{\Omega}^T)^T$.
\end{itemize}
Because the $(2,2)$ block is zero, the system is indefinite (a saddle-point problem). We use block-diagonal preconditioning:
\begin{equation}
    \mathbf{P}^{-1} = \begin{pmatrix} \tilde{\mathbf{M}}^{-1} & \mathbf{0} \\ \mathbf{0} & (\mathbf{K}^T \tilde{\mathbf{M}}^{-1} \mathbf{K})^{-1} \end{pmatrix}
\end{equation}
which clusters the eigenvalues of the preconditioned operator around $1$, allowing GMRES to converge in a small number of iterations (typically $15 - 30$ iterations).
}

% ==============================================================================
\section{Category 4: Pair Interactions, Drafting \& Far-Field Asymptotics}
% ==============================================================================

\question{Q9}{Why do two identical spheres settling side-by-side fall faster than an isolated sphere ($V_z / V_0 > 1$)? Where does the theoretical asymptotic factor $1 + \frac{3}{4}\frac{R}{d}$ come from?}

\shortanswer{Each settling sphere entrains fluid in its downward viscous wake. In Stokes flow, this wake velocity decays as the Green's function $u_z \approx \frac{F}{6\pi\eta} \frac{3}{4d}$. When a second sphere is placed at separation $d$, it is advected downward by the primary sphere's wake, leading to cooperative acceleration: $V_z / V_0 \approx 1 + \frac{3}{4}(R/d)$.}

\deepdetails{
For an isolated sphere, Stokes drag is $F = 6\pi\eta R V_0 \implies V_0 = F / (6\pi\eta R)$. 
In the presence of body 1 at the origin exerting force $\mathbf{F} = (0, 0, -F)$, the fluid velocity at body 2 located at $\mathbf{d} = (d, 0, 0)$ is given by the Oseen tensor:
\begin{equation}
    \mathbf{u}(\mathbf{d}) = \frac{1}{8\pi\eta d} \left( \mathbf{I} + \frac{\mathbf{d}\mathbf{d}^T}{d^2} \right) \begin{pmatrix} 0 \\ 0 \\ -F \end{pmatrix} = -\frac{F}{8\pi\eta d} \hat{\mathbf{z}}
\end{equation}
Substituting $F = 6\pi\eta R V_0$:
\begin{equation}
    u_z(\mathbf{d}) = -\frac{6\pi\eta R V_0}{8\pi\eta d} = -\frac{3}{4} \left(\frac{R}{d}\right) V_0
\end{equation}
By Fax\'en's theorem, the velocity of sphere 2 is its isolated velocity plus the undisturbed fluid wake velocity from sphere 1:
\begin{equation}
    V_z(d) = V_0 + |u_z(\mathbf{d})| = V_0 \left( 1 + \frac{3}{4}\frac{R}{d} \right) \implies \frac{V_z(d)}{V_0} \approx 1 + \frac{3}{4}\left(\frac{R}{d}\right)
\end{equation}
Adding higher-order reflections (Stimson-Jeffery continuum theory) yields $1 + \frac{3}{4}(R/d) + \frac{9}{16}(R/d)^2 + \mathcal{O}(d^{-3})$.
}

\vspace{0.4cm}
\question{Q10}{Why do two identical spheres settling side-by-side maintain a strictly constant separation distance ($d(t) \equiv d_0$)? Why don't they attract or repel?}

\shortanswer{Due to the fore-aft and lateral reflection symmetry of the Stokes equations. Any lateral velocity component would violate time-reversal and reflection symmetry. Spheres fall along strictly parallel, straight vertical paths.}

\deepdetails{
Consider two identical spheres at $(-d/2, 0, z)$ and $(d/2, 0, z)$ falling under $\mathbf{F} = -F \hat{\mathbf{z}}$.
\begin{enumerate}
    \item Reflection across the bisecting plane $x = 0$ maps Sphere 1 to Sphere 2 and vice-versa, requiring $\dot{x}_1 = -\dot{x}_2$.
    \item In Stokes flow, reversing the direction of time ($t \to -t$) reverses forces and velocities ($\mathbf{F} \to -\mathbf{F}, \mathbf{u} \to -\mathbf{u}$).
    \item However, reversing vertical gravity and performing a $180^\circ$ rotation about the $y$-axis restores the original system geometry.
    \item If the spheres attracted during downward fall, they would also have to attract during upward motion, but time-reversal requires them to repel. The only value consistent with both symmetries is $\dot{x}_1 = \dot{x}_2 = 0$.
\end{enumerate}
In our simulation data (`sphere_pair_trajectories.png`, Panel C), inter-body distance $d(t)/R$ is strictly constant ($0.000\%$ lateral drift).
}

% ==============================================================================
\section{Category 5: The Crowley Instability in Falling Arrays}
% ==============================================================================

\question{Q11}{What is the Crowley Instability (1971), and how did your 5-body linear horizontal chain simulations validate it?}

\shortanswer{J.M. Crowley (1971) proved that an initially horizontal array of sedimenting particles is unstable to spatial perturbations. Closer particles fall faster due to mutual drafting; by continuity, fluid in-rush laterally pulls them closer, triggering runaway sagging into an inverted-U chevron. Our 5-body simulations directly proved this: the center particle falls $15\% - 37\%$ faster than the edges, forming the exact sagging profile that precedes droplet pinch-off.}

\deepdetails{
In a linear array of $M$ particles:
\begin{itemize}
    \item Edge particles ($k=1, M$) have neighbors on only one side;
    \item The central particle ($k = \lceil M/2 \rceil$) experiences additive downward wakes from both left and right flanks.
\end{itemize}
From our simulation data across separation ratios $d/R \in [2.0, 10.0]$:
\begin{itemize}
    \item \textbf{Spheres ($N = 42$):} Center velocity exceeds edges by $+15.7\%$ at $d/R = 2.0$, peaking at $+27.5\%$ at $d/R = 3.0$ ($V_{\text{center}} / V_{\text{edge}} = 1.2754$).
    \item \textbf{Discs ($N = 37$):} Peaking at $+36.2\%$ center speedup at $d/R = 3.0$ ($V_{\text{center}} / V_{\text{edge}} = 1.3622$).
    \item \textbf{Cylinders ($N = 86$):} Peaking at $+37.2\%$ center speedup at $d/R = 3.0$ ($V_{\text{center}} / V_{\text{edge}} = 1.3721$).
    \item \textbf{Prolate Ellipsoids ($N = 42$):} Peaking at $+24.8\%$ at $d/R = 4.0$ ($V_{\text{center}} / V_{\text{edge}} = 1.2483$).
\end{itemize}
This persistent velocity gradient produces the inverted sagging chevron visible in Panels A and B of our 5-body dynamics figures (`figures/*_5body_dynamics.png`).
}

\vspace{0.4cm}
\question{Q12}{In your 5-body pentagonal ring cluster, why did all 5 particles sediment at the exact same velocity, and why was the speedup so massive ($V_z / V_0 \approx 2.44 - 2.50$)?}

\shortanswer{By $C_5$ five-fold discrete rotational symmetry, every particle in the regular pentagon occupies an identical hydrodynamic environment, enforcing identical sedimentation speeds. The massive $+144\% - 150\%$ speedup arises from constructive superposition: each particle sits in the combined drafting wake of four neighboring bodies.}

\deepdetails{
For a regular polygon of $M$ bodies on a circle of chord distance $d$, each particle $i$ experiences wake advection from all $j \ne i$:
\begin{equation}
    V_z(M) = V_0 \left[ 1 + \frac{3}{4}\left(\frac{R}{d}\right) \sum_{j \ne i} \frac{d}{d_{ij}} \right]
\end{equation}
For a pentagon ($M = 5$), the chord distances to the other four vertices are $d, d, \phi d, \phi d$, where $\phi = \frac{1+\sqrt{5}}{2} \approx 1.618$ is the golden ratio.
\begin{equation}
    \sum_{j \ne i} \frac{d}{d_{ij}} = 1 + 1 + \frac{1}{1.618} + \frac{1}{1.618} = 2 + 1.236 = 3.236
\end{equation}
Therefore:
\begin{equation}
    \frac{V_z(M=5)}{V_0} \approx 1 + \frac{3}{4} (3.236) \left(\frac{R}{d}\right) \approx 1 + 2.427 \left(\frac{R}{d}\right)
\end{equation}
At $d/R = 2.0$, this theory predicts $V_z / V_0 \approx 1 + 2.427/2 = 2.21$. Our multiblob simulation yielded $V_z / V_0 = 2.4403$ for spheres and $2.4968$ for discs, with the slight excess reflecting near-field lubrication corrections.
}

% ==============================================================================
\section{Category 6: Rahul Chajwa's Work \& Keplerian Orbits in Stokes Flow}
% ==============================================================================

\question{Q13}{Who is Rahul Chajwa, and how can conservative ``Keplerian orbits'' exist in an overdamped, zero-Reynolds-number viscous fluid?}

\shortanswer{Rahul Chajwa (Stanford / ICTS-TIFR) discovered that settling anisotropic discs execute bound periodic Keplerian orbits (PRL 2019). While spheres maintain fixed spacing, discs couple tilt orientation to lateral drift. The hydrodynamic torque exerted by one disc tilts the other. In relative phase space, tilt $\theta$ acts as conjugate momentum $p$ and separation $r$ as coordinate, yielding an effective Hamiltonian $\mathcal{H}(r, p) = p^2/(2m_{\text{eff}}) + \mathcal{V}_{\text{eff}}(r)$.}

\deepdetails{
In Chajwa, Menon, and Ramaswamy (\emph{Phys. Rev. Lett.} 122, 224501, 2019):
\begin{enumerate}
    \item \textbf{Drift from Tilt:} An inclined disc with normal vector tilted by angle $\theta$ experiences anisotropic drag ($C_\perp > C_\parallel$). Sedimenting along $-z$ induces horizontal drift: $\dot{r} \propto \sin(2\theta) \approx 2\theta$.
    \item \textbf{Torque from Wake:} The shear flow $\partial u_z / \partial r$ in the wake of the partner disc exerts a hydrodynamic torque $\mathcal{T}_H \propto -1/r^2$.
    \item \textbf{Conjugate Mapping:} Setting $p \equiv \theta$, we have $\dot{r} = p / m_{\text{eff}}$, and $\dot{p} = \dot{\theta} = \mathcal{T}_H = -\partial \mathcal{V}_{\text{eff}} / \partial r$.
    \item \textbf{Conservative Orbits:} This forms an exact canonical Hamiltonian system:
    \begin{equation}
        \mathcal{H} = \frac{p^2}{2 m_{\text{eff}}} + \mathcal{V}_{\text{eff}}(r) = E = \text{constant}
    \end{equation}
    For bound states ($E < 0$), the two discs continuously oscillate in separation $r(t)$ and tilt $\theta(t)$ without touching, tracing Keplerian ellipses in phase space!
\end{enumerate}
}

\pitfall{Do not say ``inertia causes the orbit''. The fluid has zero inertia ($Re = 0$). The ``momentum'' is an \emph{effective geometric momentum} arising from the orientation degree of freedom.}

\vspace{0.4cm}
\question{Q14}{In Chajwa's PRX 2020 paper, how do disk arrays suppress the Crowley clumping instability?}

\shortanswer{In a chain of spheres, particles have no internal rotational freedom, so clumping is unconditional. In a chain of discs, orientation couples to translation, acting like an emergent mass-and-spring lattice. This emergent elasticity supports propagating acoustic-like kinematic wave modes (phonons) that disperse disturbances and convert exponential clumping into stable waves or algebraic growth.}

\deepdetails{
In Chajwa et al. (\emph{Phys. Rev. X} 10, 041016, 2020), perturbing a 1D disc lattice reveals the dispersion relation $\omega(k)$ for perturbation wavenumber $k$. For spherical chains, the growth rate $\sigma(k) = \text{Im}[\omega] > 0$ for all wavenumbers, leading to Crowley instability. For discs, coupling between position displacement $\delta x$ and tilt angle $\delta \theta$ introduces a real propagating frequency $\text{Re}[\omega] \ne 0$ (kinematic wave modes), stabilizing the array across specific lattice orientations.
}

% ==============================================================================
\section{Category 7: Industrial & Environmental Applications}
% ==============================================================================

\question{Q15}{Industrial inkjet printing involves high ejection speeds, viscoelastic inks, and surface tension. Isn't your Stokes sedimentation model completely irrelevant to inkjet technology?}

\shortanswer{No. While droplet ejection at the nozzle involves high shear and inertia, the subsequent multi-droplet flight and reservoir stability are fundamentally governed by viscous wake interactions. Idealized Stokes sedimentation isolates the pure hydrodynamic drafting mechanism, providing the exact baseline required to model tandem droplet capture, satellite merging, and colloidal aggregation in micro-channels.}

\deepdetails{
In drop-on-demand printing:
\begin{itemize}
    \item Droplets of diameter $D \approx 20\,\mu\text{m}$ traveling in tandem at spacing $d \sim 2 - 5 D$;
    \item The leading droplet creates an elongated wake where the trailing droplet experiences reduced aerodynamic/viscous drag;
    \item Trailing droplet velocity increases by $\Delta V / V_0 \approx \frac{3}{4}(D/d)$, exactly mirroring our two-body drafting curves;
    \item Inside ink conduits, pigment particles ($d_p \sim 100\,\text{nm}$) sediment under gravity at $Re \sim 10^{-6}$, where multi-body clustering causes nozzle clogging.
\end{itemize}
Stokesian hydrodynamics provides the clean, multiphysics-free ground truth for these wake interactions.
}

\vspace{0.4cm}
\question{Q16}{What is marine snow, and how does sedimentation dynamics impact global ocean carbon models?}

\shortanswer{Marine snow consists of macroscopic organic aggregates sinking from the ocean surface to the seafloor, sequestering gigatons of carbon (the biological carbon pump). As demonstrated by Chajwa and Prakash (Science 2024), settling aggregates produce trailing solute wakes (``comet tails'') and multi-body drafting that accelerates sinking velocities, directly altering carbon sequestration timescales.}

% ==============================================================================
\section{Category 8: Orientation Coupling & Literature Benchmarking}
% ==============================================================================

\question{Q17}{Why does an inclined ellipsoid drift horizontally ($V_x \ne 0$) when falling under pure vertical gravity?}

\shortanswer{Because the hydrodynamic resistance tensor of an anisotropic particle is non-isotropic ($C_\perp > C_\parallel$). When tilted at angle $\theta$, the drag force is not anti-parallel to velocity; resolving forces yields an off-diagonal mobility component ($M_{xz} \ne 0$) that drives lateral translation.}

\deepdetails{
Let $\hat{\mathbf{n}}$ be the orientation axis of a prolate spheroid. The drag tensor in body coordinates is diagonal with perpendicular drag $C_\perp$ and parallel drag $C_\parallel$ ($C_\perp > C_\parallel$).
Transforming to laboratory coordinates for a tilt angle $\theta$:
\begin{equation}
    \mathbf{K}_{\text{drag}} = C_\parallel \hat{\mathbf{n}}\hat{\mathbf{n}}^T + C_\perp (\mathbf{I} - \hat{\mathbf{n}}\hat{\mathbf{n}}^T)
\end{equation}
Under purely vertical force $\mathbf{F} = (0, 0, -F)^T$, the resulting velocity is $\mathbf{V} = \mathbf{K}_{\text{drag}}^{-1} \mathbf{F}$:
\begin{equation}
    V_x = \frac{1}{2} \left( \frac{1}{C_\parallel} - \frac{1}{C_\perp} \right) F \sin(2\theta)
\end{equation}
\begin{itemize}
    \item At $\theta = 0^\circ$ and $\theta = 90^\circ$: $\sin(2\theta) = 0 \implies V_x = 0$ (pure vertical fall);
    \item At $\theta = 45^\circ$: $\sin(2\theta) = 1 \implies$ **Maximum lateral drift velocity**.
\end{itemize}
This was directly confirmed in our simulation script `plot_angle_drift.py` and replicated Figure 3 from *Phys. Rev. E 109, 065302 (2024)*.
}

% ==============================================================================
\section{Category 9: Machine Learning Surrogate Extrapolation}
% ==============================================================================

\question{Q18}{Why train Machine Learning models (Random Forest / XGBoost) when your multiblob code already solves the exact Stokes equations?}

\shortanswer{Computational complexity. Multiblob GMRES solves scale as $\mathcal{O}(N^2)$ to $\mathcal{O}(N^3)$, taking seconds to minutes per shape. For high-throughput industrial screening across millions of arbitrary particulate shapes, numerical simulation is too slow. The ML surrogate model trains on our high-fidelity multiblob data and predicts terminal velocity and drafting in sub-milliseconds ($< 4\,\text{ms}$, $> 500\times$ faster) with $R^2 > 0.99$.}

\deepdetails{
Profiling from our workspace:
\begin{itemize}
    \item $N = 12 \implies 0.07\,\text{s}$ per step;
    \item $N = 162 \implies 0.52\,\text{s}$ per step;
    \item $N = 2562 \implies 63.76\,\text{s}$ per step.
\end{itemize}
Our collaborator extracts 27 geometric features (volume, surface area, sphericity $\Psi$, aspect ratios, moments of inertia, chirality) from 3D surface meshes and trains Random Forest and XGBoost regressors.
On a test suite of 263 verified multiblob simulations:
\begin{itemize}
    \item Terminal velocity $V_z$: Test $R^2 = 0.9958$, RMSE $= 5.92\,\text{units/s}$ ($< 3\%$ error);
    \item Lateral drift $V_x$: Test $R^2 = 0.9593$;
    \item Inference time: $2.8 - 4.0\,\text{ms}$ ($> 500\times$ speedup).
\end{itemize}
}

\vspace{0.4cm}
\question{Q19}{What is Brenner's Theorem for regular polyhedra (such as cubes), and how did you use it to validate the ML model?}

\shortanswer{Howard Brenner (1963) proved that any particle possessing tetrahedral, octahedral, or icosahedral cubic point-group symmetry has an isotropic hydrodynamic resistance tensor: its terminal settling velocity is completely invariant to orientation ($\theta \in [0^\circ, 90^\circ]$) and lateral drift is identically zero ($V_x \equiv 0$). We tested this on regular multiblob cubes ($N=98$) and confirmed that our simulation and ML surrogate preserved this invariant behavior.}

% ==============================================================================
\section{Category 10: Classic Defense Trap Questions}
% ==============================================================================

\question{Q20}{Did you encounter ``Stokes' Paradox'' in your cylinder simulations? Why or why not?}

\shortanswer{No. Stokes' paradox states that there is no steady, inertialess 2D Stokes flow past an infinitely long cylinder in an unbounded fluid. However, our cylinders are **three-dimensional and finite in length** (aspect ratio $L / 2R_{\text{cyl}} \approx 6.5$). In 3D, fluid can flow around the ends of the cylinder, allowing the 3D Green's function to decay as $1/r$ and avoiding any mathematical paradox.}

\deepdetails{
In 2D Stokes flow, the 2D Green's function decays logarithmically ($\ln r$), which cannot satisfy the uniform far-field boundary condition $\mathbf{u} \to \mathbf{U}_\infty$ as $r \to \infty$ (Stokes' Paradox, resolved by Oseen's inertia correction). In 3D, the Stokeslet decays as $1/r \to 0$ as $r \to \infty$. Because our multiblob cylinders have finite ends ($N = 86$), the 3D Stokes equations possess a unique, well-behaved solution.
}

\vspace{0.4cm}
\question{Q21}{If gravity is a conservative potential force and Stokes flow is purely dissipative, where does the potential energy of the sedimenting particle go?}

\shortanswer{The lost gravitational potential energy is converted entirely into heat via viscous dissipation in the fluid: $\frac{dE_p}{dt} = \mathbf{F}_{\text{ext}} \cdot \mathbf{V} = -\int_{\Omega} 2\eta (\mathbf{E}:\mathbf{E}) \, dV < 0$, where $\mathbf{E} = \frac{1}{2}(\nabla \mathbf{u} + (\nabla \mathbf{u})^T)$ is the fluid rate-of-strain tensor.}

\vspace{0.4cm}
\question{Q22}{In a suspension of billions of particles, why can't we simply integrate the Oseen tensor $1/r$ over the entire fluid volume? What is the ``sedimentation divergence paradox'', and how did Batchelor resolve it?}

\shortanswer{Integrating $1/r$ over a 3D volume produces a conditionally convergent or divergent integral ($\int_0^\infty \frac{1}{r} r^2 dr \to \infty$). This is the classical sedimentation paradox. G.K. Batchelor (1972) resolved this by recognizing that in a closed container, the downward flux of particles must be balanced by an upward fluid backflow ($\langle \mathbf{u} \rangle = \mathbf{0}$), subtracting the mean volume flow and yielding the convergent settling formula: $V(\phi) = V_0 (1 - 6.55 \phi)$.}

% ==============================================================================
\section{Category 11: RigidMultiblobsWall Software Architecture \& Implementation}
% ==============================================================================

\question{Q23}{What software framework powers your simulations? Who authored it, and what are its key modules?}

\shortanswer{We use \texttt{RigidMultiblobsWall}, an open-source high-performance Stokesian dynamics framework authored by Balboa Usabiaga, Delong, Delmotte, Sprinkle, and Donev (Courant Institute / Flatiron Institute / CAMCoS 2016, JCP 2017). Its core modules include \texttt{mobility/} for RPY and Blake Green's function solves, \texttt{multi\_bodies/} for dynamic and quasi-steady multi-body integration, and \texttt{quaternion\_integrator/} for rotational tracking.}

\deepdetails{
The repository architecture consists of:
\begin{itemize}
    \item \texttt{mobility/mobility.py}: Evaluates the RPY mobility tensor $\mathbf{M}$, the Rotne-Prager-Blake (RPB) tensor for walls, and the block-diagonal preconditioned GMRES saddle-point solver.
    \item \texttt{multi\_bodies/multi\_bodies.py}: Time-dependent integrator managing deterministic forward Euler, Adams-Bashforth, and stochastic RFD/Euler-Maruyama solvers.
    \item \texttt{multi\_bodies/multi\_bodies\_utilities.py}: Quasi-steady solver for single-step mobility solves (\texttt{scheme mobility}), drag resistance (\texttt{scheme resistance}), and $6 \times 6$ rigid-body mobility tensors (\texttt{scheme body\_mobility}).
    \item \texttt{quaternion\_integrator/quaternion\_integrator.py}: Integrates rigid body orientations on the $S^3$ unit hypersphere with norm renormalization.
    \item \texttt{multi\_bodies/forces.cpp}: Optional C++/Numba steric repulsion routines for inter-particle Debye/Yukawa potentials.
\end{itemize}
}

\vspace{0.4cm}
\question{Q24}{What are \texttt{*.vertex} and \texttt{*.clones} files, and how does the framework represent arbitrary rigid bodies?}

\shortanswer{A \texttt{*.vertex} file defines the body's internal geometry by listing the Cartesian coordinates $(x_i, y_i, z_i)$ of all constituent blobs relative to its center of mass. A \texttt{*.clones} file defines the global state of each rigid body instance in the simulation, listing its global location $(x, y, z)$ and orientation unit quaternion $(q_0, q_1, q_2, q_3)$.}

\deepdetails{
For an arbitrary shape (such as a disc, cylinder, or trefoil knot):
\begin{enumerate}
    \item The continuous geometry is discretized into $N$ blobs of hydrodynamic radius $a$ and written to \texttt{Structures/shape\_N.vertex}.
    \item Multiple bodies (e.g., pairs at separation $d$, 5-body chains, or ring clusters) are placed in the simulation by defining their positions and orientations in \texttt{Structures/shape\_M\_bodies.clones}.
    \item Blobs move according to rigid body kinematics $\mathbf{u}_i = \mathbf{V} + \mathbf{\Omega} \times (\mathbf{x}_i - \mathbf{q})$. The framework assembles the kinematic matrix $\mathbf{K}$ and solves for the unknown constraint forces $\boldsymbol{\lambda}$ and rigid body velocities $\mathbf{U} = (\mathbf{V}^T, \mathbf{\Omega}^T)^T$.
\end{enumerate}
}

\vspace{0.4cm}
\question{Q25}{How does the code switch between unbounded free-space and wall-bounded domains, and what mathematical kernel is activated?}

\shortanswer{By setting the \texttt{domain} parameter in \texttt{input.dat}. In \texttt{domain free\_space}, the code uses the free-space RPY tensor based on the Oseen Stokeslet. In \texttt{domain one\_wall}, the code activates the Rotne-Prager-Blake (RPB) tensor, which uses Blake's image system (Stokeslet + inverted image Stokeslet + Stokes-dipole + potential doublet) to strictly enforce no-slip $\mathbf{u} = \mathbf{0}$ at $z = 0$.}

\vspace{0.4cm}
\question{Q26}{What output files are generated, and how did you verify numerical convergence and reproducibility?}

\shortanswer{The framework outputs \texttt{*.velocities.dat} ($6$-DOF velocities), \texttt{*.clones} (trajectory positions/quaternions), \texttt{*.time} (execution wall-clock seconds), and \texttt{*.info} (GMRES iteration counts and residual logs). We verified solver convergence by setting GMRES tolerance to $10^{-8}$, confirming that velocity residuals were below machine noise and results were fully deterministic via \texttt{*.random\_state}.}

\vspace{0.4cm}
\question{Q27}{How does your automated simulation workflow wrap around \texttt{RigidMultiblobsWall}?}

\shortanswer{We built an automated Python orchestration layer (\texttt{custom\_simulations/}) that programmatically generates parameterized \texttt{*.vertex} and \texttt{*.clones} configurations across separation ratios $d/R \in [2.0, 12.0]$, constructs isolated run directories with dedicated \texttt{input.dat} parameters, executes the \texttt{RigidMultiblobsWall} solvers via subprocess, extracts telemetry into standardized scientific whitespace-delimited \texttt{.dat} tables, and renders publication-ready 4-panel trajectory plots.}

\end{document}

